\subsection{Работа программного средства}

К основному функционалу программного средства можно отнести
следующие компоненты:

\begin{itemize}
    \item получение списка директорий;
    \item получение списка файлов в директории;
    \item открытие файла;
\end{itemize}

\subsubsection{Получение списка директорий}

Получение списка директорий начинается с проверки наличия
placeholder-элементов в узле дерева
каталогов и их удаления.

После удаления placeholder-элементов функция проверяет
данные узла на существование и проверяет, был ли узел
уже расширен ранее, чтобы избежать повторной загрузки
подкаталогов.

Если узел ещё не расширен, формируется паттерн поиска
для директории и открывается поиск файлов через функцию
FindFirstFileW. При успешном открытии поиска выполняется
перебор всех найденных элементов: пропускаются файлы,
пропускаются системные директории, для каждой найденной
директории проверяется наличие подкаталогов и элемент
добавляется в дерево с соответствующей пометкой о наличии
дочерних элементов.

\begin{lstlisting}[style=CodeListing]
void ExpandTreeItem(HWND hTree, HTREEITEM hItem) {
    // Remove placeholder children
    HTREEITEM hChild = TreeView_GetChild(hTree, hItem);
    while (hChild) {
        TVITEMW tvi{}; tvi.mask = TVIF_TEXT; tvi.hItem = hChild;
        wchar_t buf[8]; tvi.pszText = buf; tvi.cchTextMax = 8;
        TreeView_GetItem(hTree, &tvi);
        if (wcscmp(buf, L"...") == 0) {
            TreeView_DeleteItem(hTree, hChild);
            break;
        }
        hChild = TreeView_GetNextSibling(hTree, hChild);
    }

    TVITEMW tvi{}; tvi.mask = TVIF_PARAM; tvi.hItem = hItem; 
    TreeView_GetItem(hTree, &tvi);
    NodeData *data = reinterpret_cast<NodeData*>(tvi.lParam);
    if (!data || data->expanded) return;

    std::wstring pattern = Utils::JoinPath(data->path, L"*");
    WIN32_FIND_DATAW fd{};
    HANDLE h = FindFirstFileW(pattern.c_str(), &fd);
    if (h == INVALID_HANDLE_VALUE) return;
    do {
        if ((fd.dwFileAttributes & FILE_ATTRIBUTE_DIRECTORY) == 0) continue;
        std::wstring name = fd.cFileName;
        if (name == L"." || name == L"..") continue;
        std::wstring full = Utils::JoinPath(data->path, name);
        bool hasChildren = HasSubdirs(full);
        InsertTreeItem(hTree, hItem, name, full, hasChildren);
    } while (FindNextFileW(h, &fd));
    FindClose(h);
    data->expanded = true;
}
\end{lstlisting}

После завершения перебора поиск закрывается, и узел
помечается как расширенный, чтобы предотвратить повторную
загрузку при следующем раскрытии.

Если в процессе работы алгоритма произошла ошибка, функция
завершается без добавления элементов, и узел остаётся в
исходном состоянии.

\subsubsection{Получение списка файлов в директории}

Получение списка файлов в директории начинается с очистки
списка интерфейса и вектора элементов. Далее выполняется
нормализация пути директории для разрешения символических
ссылок и коротких путей.

После нормализации пути проверяется существование директории
через GetFileAttributesW. Если путь не существует,
отображается соответствующее сообщение об ошибке, и функция
завершается.

Если путь существует, проверяется, что это действительно
директория, а не файл. Если переданный путь указывает на
файл, отображается сообщение "Файл пуст", и функция
завершается.

\begin{lstlisting}[style=CodeListing]
int idx = 0;
do {
    std::wstring name = fd.cFileName;
    if (name == L"." || name == L"..") continue;
    bool isDir = (fd.dwFileAttributes & FILE_ATTRIBUTE_DIRECTORY) != 0;

    ListItem ri{ name, Utils::JoinPath(normalizedPath, name), isDir };
    items.push_back(ri);

    LVITEMW item{};
    SHFILEINFOW sfi{};
    UINT attrs = isDir ? FILE_ATTRIBUTE_DIRECTORY : FILE_ATTRIBUTE_NORMAL;
    SHGetFileInfoW(ri.fullPath.c_str(), attrs, &sfi, sizeof(sfi), 
                  SHGFI_SYSICONINDEX | SHGFI_USEFILEATTRIBUTES | SHGFI_SMALLICON);

    item.mask = LVIF_TEXT | LVIF_PARAM | LVIF_IMAGE;
    item.iItem = idx;
    item.pszText = const_cast<wchar_t*>(name.c_str());
    item.lParam = static_cast<LPARAM>(idx);
    item.iImage = sfi.iIcon;
    ListView_InsertItem(hList, &item);

    std::wstring typeText = isDir ? L"Dir" : L"File";
    ListView_SetItemText(hList, idx, 1, const_cast<wchar_t*>(typeText.c_str()));

    std::wstring mod = Utils::FileTimeToWString(fd.ftLastWriteTime);
    ListView_SetItemText(hList, idx, 2, const_cast<wchar_t*>(mod.c_str()));
    ++idx;
} while (FindNextFileW(h, &fd));
FindClose(h);
\end{lstlisting}

Для директории формируется паттерн поиска и открывается
поиск файлов. При успешном открытии выполняется перебор
всех найденных элементов: пропускаются системные директории,
для каждого элемента определяется его тип, создаётся объект
элемента списка, получается системная иконка, и элемент
добавляется в список интерфейса с установкой текста колонок
"Type" и "Modified".

После завершения перебора поиск закрывается, и список файлов
отображается пользователю.

\subsubsection{Открытие файла}

Открытие файла начинается с обработки события двойного клика
пользователя по элементу списка. Функция получает индекс
выбранного элемента и проверяет его корректность.

После получения индекса функция извлекает соответствующий
элемент из списка элементов директории и проверяет его тип.
Если выбранный элемент является директорией, выполняется
переход в эту директорию через функцию NavigateTo.

\begin{lstlisting}[style=CodeListing]
int idx = act->iItem;
if (idx >= 0 && idx < (int)state->listItems.size()) {
    const Filesystem::ListItem &ri = state->listItems[idx];
    if (ri.isDir) {
        NavigateTo(hwnd, *controls, *state, ri.fullPath, true);
    } else {
        HINSTANCE h = ShellExecuteW(hwnd, L"open", ri.fullPath.c_str(), 
                                   nullptr, nullptr, SW_SHOWNORMAL);
        if ((INT_PTR)h <= 32) {
            MessageBoxW(hwnd, L"Не удалось открыть файл", L"Ошибка", 
                       MB_ICONERROR | MB_OK);
        }
    }
}
\end{lstlisting}

Если выбранный элемент является файлом, выполняется попытка
открыть файл через функцию ShellExecuteW с параметром "open",
который передаёт управление файлом операционной системе для
выбора соответствующего приложения.

При успешном открытии файла функция ShellExecuteW возвращает
значение больше 32. Если возвращаемое значение меньше или
равно 32, это означает ошибку открытия файла, и пользователю
отображается сообщение об ошибке "Не удалось открыть файл".

Если файл успешно открыт, операционная система запускает
соответствующее приложение для работы с файлом, и
пользователь может работать с содержимым файла.